\chapter{Evaluation}
\label{ch:evaluation}

This chapter evaluates the program as it stands after commit \commit{448d4f1}. We
consider what it does well, where it falls short, and what could be done next.

\section{Strengths}
\label{sec:strengths}

\textbf{The movement maths is visible and checkable.} Every distance and angle is
converted with an equation that can be worked out by hand (\eref{eq:drive},
\eref{eq:turn}). When a move is wrong, we can check each number in turn instead of
guessing what a hidden library function is doing.

\textbf{Calibration happens on the robot.} \code{turn_tuner()} lets the driver tune
\code{TURN_SCALE} with the controller in a few minutes, without a laptop. This is
useful at a competition, where the floor surface may differ from the practice
field and turns need retuning.

\textbf{Faults can be found without a laptop.} The vision report, the port scanner,
and the claw diagnostic all draw on the Brain screen. These tools were what solved
the silent vision failure (\sref{sec:ch-silent}).

\textbf{The driver is always in control.} The cone assist can be cancelled by any
button or stick, and the claw diagnostic cannot crash the drive loop. A failure in
an optional feature therefore never takes control away from the driver.

\textbf{Settings are in one place, and the reasoning is written down.} Every
tunable number is in \code{config.py}, and \code{README.md} and \code{AGENTS.md}
explain the decisions and the pitfalls. Someone new to the code can understand it
without repeating our mistakes.

\section{Limitations}
\label{sec:limitations}

\subsection{Open-Loop Turns}

The autonomous routine is \emph{open-loop}: the robot turns the wheels a calculated
number of degrees, but never checks which way it is actually facing. If a wheel
slips, the error goes unnoticed and affects the rest of the path. \code{TURN_SCALE}
corrects the average error, but not one-off slips. An Inertial sensor, which
measures rotation directly, would allow \emph{closed-loop} turns that keep turning
until the robot really faces the right way. The SDK supports this, but we have not
confirmed whether the sensor is fitted to the robot.

\subsection{No Competition Template}

\code{main()} runs one mode chosen by hand (\sref{sec:modes}). At an official match,
the field control system needs the program to register both
\code{autonomous()} and \code{user_control()} with a \code{Competition} object, so that
it can start each one at the right time. The program does not do this yet, so
\code{main()} must be changed before the robot is used in an official match.

\subsection{An Inactive Safeguard}

\code{fit_turn()} only acts when forward power plus turn power exceeds 100
(\sref{sec:taper}). With the current settings, this never happens in normal
driving: forward power is at most 50\% (\code{SPEED_NORMAL}), and turn power is at
most $50 \times 0.8 = 40\%$, further reduced by the taper. Holding R1 can exceed
100, but R1 skips \code{fit_turn()}. The function is therefore inactive. It is
harmless, but it suggests the speed settings were changed after it was written,
and it should either be removed or kept deliberately for if the speeds are raised.

\subsection{Wide Colour Tolerances}

The colour descriptions use wide hue and saturation tolerances
(\sref{sec:describe}). This helps detection in poor lighting, but increases the
chance of matching something that is not the cone. The colour code and the
colour-pair distance check reduce this risk, but do not remove it.

\subsection{Wall Squaring Assumes a Wall}

\code{square_to_wall()} assumes the robot starts with its back near a flat wall. If
it starts too far away, it reverses for the full 2.5~s and then begins the path
from wherever it stopped, with no warning.

\subsection{Documentation Out of Date}

The README and AGENTS documents still say the program downloads to slot 2, but
commit \commit{76b036f} moved it to slot 3. This is minor, but it shows that
documentation must be updated together with the code.

\section{Future Work}

\begin{enumerate}
  \item Add a \code{Competition} object so the program works under field control.
  \item Use the Inertial sensor, if fitted, for closed-loop turns.
  \item Show a warning on the Brain screen if \code{square_to_wall()} times out
    without finding a wall.
  \item Narrow the colour tolerances once the competition lighting is known.
\end{enumerate}
